% VI34-DETAILED-PROBLEM-19
\begin{problem}[VI.34.P19: Finite graded presentations]\label{prob:vi34-19}
Let \(M\) be generated by homogeneous elements \(m_i\) of degrees \(d_i\). Construct a graded free presentation by
shifts, assuming the relation module is finitely generated by homogeneous elements.

\textbf{Provenance.} Canonical VI/34 extension.
\end{problem}

\ifdefined\FullProblemDossiers
\begin{solution}
For each homogeneous generator \(m_i\in M_{d_i}\), the element \(1\in S(-d_i)\) has shifted degree \(d_i\).
Therefore the assignment
\[
1_i\longmapsto m_i
\]
defines a degree-zero graded surjection
\[
F_0:=\bigoplus_iS(-d_i)\twoheadrightarrow M.
\]
Let
\[
K=\ker(F_0\to M).
\]
Because the map has degree zero, \(K\) is a graded submodule.

Assume \(K\) is generated by homogeneous relations \(r_j\) of degrees \(e_j\). Then each relation defines a
degree-zero map from \(S(-e_j)\), and together they give
\[
F_1:=\bigoplus_jS(-e_j)\longrightarrow F_0
\]
with image \(K\). Thus
\[
\boxed{
F_1\longrightarrow F_0\longrightarrow M\longrightarrow0
}
\]
is a graded presentation by finite direct sums of shifts when both generator and relation sets are finite.

If \(S\) is Noetherian and \(M\) is finitely generated graded, the kernel \(K\) is finitely generated, and its
homogeneous components can be chosen as homogeneous generators.
\end{solution}
\fi
